\hspace{3mm}

\centerline{\bf \Huge Графы}

\centerline{\bf \large Базовый уровень}

\begin{flushright}
        \scriptsize{}
\end{flushright}

\problem{0.1} а) Сытый марсианский кот Васька поймал 6 марсианских мышек (у марсианских мышек --- 3 хвоста) и стал связывать их хвостами: берёт два свободных хвоста, и связывает их узлом. Через некоторое время свободных хвостов не осталось. Сколько узелков ему пришлось завязать? б) На следующий день Васька поймал еще одну мышку и решил, развязав некоторые из узелков, связать эту мышку со всеми остальными. Сможет ли он это сделать так, чтобы по-прежнему не было свободных хвостов?

\textbf{Лемма о рукопожатиях}. В любом графе сумма степеней вершин --- чётное число. То есть, количество вершин, имеющих нечетную степень, --- чётно. Эта лемма справедлива также и для каждой компоненты связности графа в отдельности.

\problem{1}
Вероника натянула во дворе веревки для белья. Причем к $6$ столбам она прикрепила по $7$ веревок, к $4$ столбам --- по 6 веревок, и к $8$ --- по $2$ веревки. Сколько всего веревок она натянула?

\problem{2}
Артём нарисовал на окружности $9$ точек и каждые две точки соединил отрезком. Сколько отрезков нарисовал Артём?

\problem{3}
Существует ли граф на 8-ми вершинах, степени которого равны:
(а) 8,6,5,4,4,3,2,2;(б) 7,7,6,5,4,2,2,1;(в) 6,6,6,5,5,3,2,2?

\problem{4}
В стране 15 городов, каждый из которых соединен дорогами не менее чем с 7 другими. Докажите, что из любого города можно добраться до любого другого.

\problem{5}
 В королевстве из каждого города выходит 100 дорог, причем из каждого города можно доехать до любого другого. Одну дорогу закрыли на ремонт. Докажите, все еще от каждого города можно добраться до любого другого.


\newpage

\hspace{3mm}

\centerline{\bf \Huge Графы}

\centerline{\bf \large Продвинутый уровень}

\begin{flushright}
        \scriptsize{}
\end{flushright}

\textbf{Определение.} Граф называется \textit{двудольным}, если его вершины можно раскрасить в два цвета так, чтобы любое ребро соединяло вершины разных цветов.

\problem{0.1} В классе каждый мальчик дружит с тремя девочками, а каждая девочка $—$ с пятью мальчиками. $17$ из них любят играть в матбой, и в классе $15$ парт. Сколько всего ребят в классе?

\problem{0.2} 
Докажите, что в двудольном графе суммы степеней вершин каждого цвета совпа- дают между собой.

\problem{1}
Докажите, что в двудольном графе не бывает циклов нечётной длины.

\problem{2}
Докажите, что дерево является двудольным графом. Верно ли обратное?

\problem{3} 
$10$ школьников решали $10$ задач. Выяснилось, что каждый школьник решил ровно две задачи, и каждую задачу решило два школьника. Докажите, что можно так организовать рассказ решений, чтобы каждый школьник рассказал одну из решенных им задач, и при этом все задачи были бы рассказаны.

\problem{4}
Можно ли расставить 777 шахматных коней на доске $2024\times 2024$ так, чтобы каждый из них бил ровно 4 других коней?

\problem{5} 
В каждой строке и в каждом столбце таблицы $8\times8$ стоит ровно две фишки. Докажите, что их можно покрасить в чёрный и белый цвет так, чтобы в каждом столбце и в каждой строке стояли разноцветные фишки.